\section{研究计划与可交付成果}

\subsection{总体技术路线}

本项目按照一条清晰主线组织工作，各阶段依次推进、前序输出为后序输入：

\begin{center}
\begin{tabular}{c}
\hline
\textbf{阶段一：数据审计与预处理} \\
$\downarrow$ \\
\textbf{阶段二：探索性数据分析} \\
$\downarrow$ \\
\textbf{阶段三：能耗基准回归预测} \\
$\downarrow$ \\
\textbf{阶段四：模型解释} \\
$\downarrow$ \\
\textbf{阶段五：异常高耗能识别与建筑级汇总} \\
$\downarrow$ \\
\textbf{阶段六：改造优先级排序与稳定性检验} \\
$\downarrow$ \\
\textbf{阶段七：辅助聚类画像} \\
$\downarrow$ \\
\textbf{阶段八：可视化与结果分析} \\
\hline
\end{tabular}
\end{center}

\subsection{各阶段任务说明}

\subsubsection{阶段一：数据审计与预处理}

\begin{itemize}
    \item \textbf{数据挖掘问题类型：}数据清洗、缺失值处理、异常值检测
    \item \textbf{输入：}2015--2024年原始CSV文件（建筑—年份粒度，约34,699条记录，47个字段）
    \item \textbf{核心处理：}
    \begin{enumerate}
        \item 跨年度字段审计：检查各年份字段名称一致性、单位一致性、建筑类型命名规范性和建筑ID稳定性
        \item 数据质量评估：统计各字段缺失率、重复记录比例、数值字段解析异常
        \item 目标泄漏字段识别与剔除：排除与SiteEUI存在确定性数量关系的同期派生字段（SourceEUI、SiteEnergyUse等）
        \item 目标变量清洗：删除SiteEUI缺失、非正值和录入错误记录，保留合理的高能耗极端值
        \item 缺失值处理：数值字段优先按同建筑类型中位数填补，仍缺失则用训练集整体中位数；类别字段填充为Unknown；同步生成缺失指示变量
        \item 异常值处理：YearBuilt大于DataYear的记录标记为异常，按建筑类型中位年份填补，保证建筑年龄非负
        \item 数据集划分：按年份划分训练集（2015--2021）、验证集（2022）和测试集（2023、2024）
    \end{enumerate}
    \item \textbf{输出：}清洗后的特征矩阵与目标向量、数据质量审计报告（字段映射表、各年字段覆盖率、缺失值统计、重复记录报告）
\end{itemize}

\subsubsection{阶段二：探索性数据分析}

\begin{itemize}
    \item \textbf{数据挖掘问题类型：}描述性统计与可视化
    \item \textbf{输入：}阶段一清洗后的数据
    \item \textbf{核心处理：}
    \begin{enumerate}
        \item SiteEUI分布分析（原始尺度与对数变换后对比）
        \item 各建筑类型样本量与能耗水平对比（箱线图）
        \item 数值特征相关性分析（热力图）
        \item 建造年份与能耗关系分析（散点图与趋势线）
        \item 缺失值分布可视化
        \item 基准数据集匹配与超耗比例分布
    \end{enumerate}
    \item \textbf{输出：}EDA可视化图表集
\end{itemize}

\subsubsection{阶段三：能耗基准回归预测}

\begin{itemize}
    \item \textbf{数据挖掘问题类型：}监督学习—回归
    \item \textbf{输入：}训练集（2015--2021）、验证集（2022）、测试集（2023、2024）
    \item \textbf{核心处理：}
    \begin{enumerate}
        \item 基线构建：全体均值基线（Baseline A）和建筑类型中位数基线（Baseline B）
        \item 特征工程：构造BuildingAge、LogGFA、面积比率、混合用途标识、缺失指示变量等建筑属性特征；类别特征采用OneHotEncoder编码（禁止对建筑类型使用LabelEncoder）
        \item 模型训练与调优：按复杂度递增顺序评估ElasticNet、随机森林、XGBoost、LightGBM四种模型，采用按年份扩展窗口滚动交叉验证（2015--2017$\rightarrow$2018，\ldots，2015--2021$\rightarrow$2022），以负MAE为目标进行随机搜索调参；训练时对目标取对数缓解右偏，评估时还原为原始尺度
        \item 泛化能力检验：主实验采用时间外测试（2023、2024年）；辅助实验按建筑ID分组交叉验证以检验对未见建筑的泛化能力
        \item 设置两组对比实验：模型A不使用ENERGY STAR Score（主结果），模型B使用ENERGY STAR Score（扩展对比）
    \end{enumerate}
    \item \textbf{输出：}最优回归模型，各模型在验证集和测试集上的MAE、RMSE、$R^2$指标（按年份、建筑类型、高能耗子集、已见/未见建筑分组报告）
\end{itemize}

\subsubsection{阶段四：模型解释}

\begin{itemize}
    \item \textbf{数据挖掘问题类型：}模型可解释性分析
    \item \textbf{输入：}最佳模型、测试集数据
    \item \textbf{核心处理：}
    \begin{enumerate}
        \item 计算测试集SHAP值
        \item 生成全局特征重要性排名（条形图、beeswarm图）
        \item 绘制关键特征（BuildingAge、LogGFA）的SHAP依赖图
        \item 绘制典型高能耗与低能耗建筑的对比瀑布图
        \item 对高度相关特征进行分组解释
    \end{enumerate}
    \item \textbf{输出：}SHAP可解释性图表集；报告中统一使用``重要预测特征''表述，明确标注SHAP反映模型关联而非因果关系
\end{itemize}

\subsubsection{阶段五：异常高耗能识别与建筑级汇总}

\begin{itemize}
    \item \textbf{数据挖掘问题类型：}异常检测（基于残差分析和统计基准）
    \item \textbf{输入：}模型预测结果、原始SiteEUI、建筑类型—年份基准数据
    \item \textbf{核心处理：}
    \begin{enumerate}
        \item 第一层——预测残差异常：按建筑类型和年度计算标准化残差，标记标准化残差显著偏高或位于同类型同年度前10\%的建筑—年份记录
        \item 第二层——同类基准超标：计算实际SiteEUI超出同类型P75基准的比例，标记超耗比例显著的记录
        \item 第三层——持续性验证：检查2022--2024年中各建筑进入高耗能前25\%和残差前10\%的年数
        \item 建筑级聚合：按OSEBuildingID将多年记录汇总为建筑级结果（2024年当前状态 + 2022--2024年持续性表现），每栋建筑仅保留一行
    \end{enumerate}
    \item \textbf{输出：}异常高耗能建筑候选列表、建筑级汇总表（含当前状态和持续性指标）
\end{itemize}

\subsubsection{阶段六：改造优先级排序与稳定性检验}

\begin{itemize}
    \item \textbf{数据挖掘问题类型：}多指标综合排序评估
    \item \textbf{输入：}建筑级汇总表
    \item \textbf{核心处理：}
    \begin{enumerate}
        \item 6维度加权评分：标准化预测残差、同类基准超耗程度、持续高耗能频率、理论超额能耗量、GHG排放强度、建筑年龄
        \item 数据质量折扣：最终得分乘以数据完整率，防止缺失严重的建筑因极端值错误排前
        \item 多权重方案对比（标准/节能优先/减碳优先/大型建筑优先）
        \item 排序稳定性验证：不同方案下Top-N重合率与Jaccard相似度、Spearman与Kendall秩相关系数、Bootstrap重采样后各建筑入选频率、各建筑类型候选占比合理性
    \end{enumerate}
    \item \textbf{输出：}全量建筑优先级排名表、Top-N候选改造清单、权重敏感性分析报告、Bootstrap入选频率统计
\end{itemize}

\subsubsection{阶段七：辅助聚类画像}

\begin{itemize}
    \item \textbf{数据挖掘问题类型：}聚类分析（辅助性）
    \item \textbf{输入：}高优先级候选建筑的特征数据
    \item \textbf{核心处理：}
    \begin{enumerate}
        \item 聚类对象限定为优先级排名较高的候选建筑（非全部建筑），聚类不再负责判断``谁异常''
        \item 聚类特征仅使用连续数值变量（建筑面积对数、建筑年龄、基准超耗比例、持续高耗能频率、GHG强度、理论超额能耗量），建筑类型不参与距离计算，仅用于聚类结果的事后解释
        \item 综合肘部法则、轮廓系数、Calinski-Harabasz指数确定K值
        \item 为各聚类计算特征均值画像并命名（如``大型老旧高超耗办公楼''、``高排放高强度酒店''等）
    \end{enumerate}
    \item \textbf{输出：}聚类画像报告与PCA降维可视化图表
\end{itemize}

\subsubsection{阶段八：可视化与结果分析}

\begin{itemize}
    \item \textbf{数据挖掘问题类型：}结果呈现与综合分析
    \item \textbf{输入：}前序各阶段输出
    \item \textbf{核心处理：}汇总各阶段关键结果，生成综合分析可视化，撰写分析报告
    \item \textbf{输出：}Jupyter Notebook交互式分析报告、项目结题报告
\end{itemize}

\subsection{可交付成果}

\begin{enumerate}
    \item 完整的Python源码（模块化结构，各阶段独立可运行，通过主控脚本串联为完整流程）
    \item Jupyter Notebook交互式分析报告
    \item 训练好的模型文件
    \item 数据质量审计报告（含字段映射、各年覆盖率、缺失值统计、重复记录核查）
    \item 可视化图表集合（EDA图表、SHAP可解释性图表、聚类画像图表）
    \item 优先级排序输出（全量排名表、Top-N候选清单、敏感性分析报告）
    \item 开题报告与项目结题报告
\end{enumerate}

